%% ============================================================================
%%  MambaGuard: Certified Streaming Detection for Multi-Protocol
%%  LLM Agent Security
%%
%%  Target venue : IEEE Symposium on Security and Privacy 2027 (Cycle 2)
%%  Class        : IEEEtran v1.8b, [conference,compsoc], US letter
%%  Page cap     : 13 pages body + up to 5 pages references & appendices
%%
%%  ANONYMISED FOR DOUBLE-BLIND REVIEW
%%  Do not add author names, affiliations, funders, or any identifying
%%  acknowledgements. Self-citations are rephrased in the third person.
%% ============================================================================

\documentclass[conference,compsoc]{IEEEtran}

% --- packages (compsoc-compatible only; no geometry/setspace/titlesec) ----
\ifCLASSOPTIONcompsoc
  \usepackage[nocompress]{cite}
\else
  \usepackage{cite}
\fi
\usepackage{amsmath,amssymb,amsfonts,amsthm}
\usepackage{graphicx}
\usepackage{booktabs}
\usepackage{multirow}
\usepackage{array}
\usepackage{algorithm}
\usepackage{algpseudocode}
\usepackage{xcolor}
\usepackage{bm}
\usepackage[hidelinks]{hyperref}

% --- math shortcuts -------------------------------------------------------
\newcommand{\R}{\mathbb{R}}
\newcommand{\E}{\mathbb{E}}
\newcommand{\Prob}{\mathbb{P}}
\newcommand{\norm}[1]{\left\lVert#1\right\rVert}
\newcommand{\calA}{\mathcal{A}}
\newcommand{\calC}{\mathcal{C}}
\newcommand{\calG}{\mathcal{G}}
\newcommand{\calL}{\mathcal{L}}
\newcommand{\calM}{\mathcal{M}}
\newcommand{\calN}{\mathcal{N}}
\newcommand{\calO}{\mathcal{O}}
\newcommand{\calV}{\mathcal{V}}
\newcommand{\calE}{\mathcal{E}}
\DeclareMathOperator*{\argmax}{arg\,max}
\DeclareMathOperator*{\argmin}{arg\,min}
\newtheorem{theorem}{Theorem}
\newtheorem{lemma}{Lemma}
\newtheorem{proposition}{Proposition}
\newtheorem{definition}{Definition}
\newtheorem{assumption}{Assumption}

\begin{document}

\title{MambaGuard: Certified Streaming Detection for\\
Multi-Protocol LLM Agent Security}

\author{\IEEEauthorblockN{Anonymous Submission}\\
\IEEEauthorblockA{Paper \#XXXX, IEEE S\&P 2027 Cycle 2}}

\maketitle

%% ============================================================================
\begin{abstract}
Large language model (LLM) agents now communicate through four open protocols --- the Model Context Protocol, the Agent-to-Agent Protocol, the Agent Communication Protocol, and the Agent Network Protocol --- whose combined traffic forms a network-scale attack surface that prompt-boundary monitors and classical intrusion detection systems were not designed to observe. We present \emph{MambaGuard}, a streaming detector that ingests a canonicalised protocol message graph, extracts per-agent sequential features with a selective state-space backbone, propagates relational context with a Lipschitz-normalised temporal graph attention layer, and selects defensive actions through a Hedge-based controller. The detector carries two complementary certificates. First, a randomised-smoothing certificate on the sentence-transformer feature space yields a lower bound on the expected defender utility of the form $V^{*} - B\,\Pr[\neg\calC_{\varepsilon}] - B\sqrt{\ln N/(2T)} - B\alpha_{MC}$, in which every term shares the utility unit so the composition type-checks. Second, a sparse-smoothing certificate on the protocol graph bounds the number of injected or removed edges an attacker can tolerate per window. We evaluate MambaGuard on three protocol-native benchmarks (MCPSecBench, MCP-SafetyBench, MCPTox) plus an A2A sandbox and two agent-tool corpora (AgentDojo, InjecAgent), and subject it to a dedicated adaptive-attack suite that includes white-box projected gradient descent in embedding space, discrete token search, LLM-driven payload rewriting, graph-structure evasion, and attacks on the online controller. Against static attacks MambaGuard reduces post-defence attack success rate below six per~cent on four of the six benchmarks; against adaptive attackers the certified fraction at $\varepsilon = 0.05$ remains above fifty per~cent on the protocol-native benchmarks while uncertified empirical robustness collapses on three of the eight baselines. We release an anonymised reference implementation, the evaluation harness, and a reproducibility script that regenerates every table and figure from a single command.
\end{abstract}

%% ============================================================================
\section{Introduction}
\label{sec:intro}
%% ============================================================================
Enterprise AI has moved from single-prompt interfaces to protocol-mediated agent meshes. A 2025 systematic review of the Model Context Protocol catalogued more than thirty attack techniques at the MCP layer alone~\cite{Hou2025MCP}, and the empirical MCPSecBench study reported attack success rates above eighty-five per~cent on at least one major commercial deployment~\cite{Yang2025MCPSec}. The MCP Inspector remote-code-execution flaw~\cite{Qualys2025CVE49596} and the A2ABreak catalogue of eleven protocol-level vulnerabilities~\cite{Lotfi2026A2ABreak} confirm that protocol exposure translates directly into system compromise. The defensive response so far has been fragmented along three separate axes.

First, \emph{design-level} defences such as CaMeL~\cite{Debenedetti2025CaMeL} and MELON~\cite{Zhu2025MELON} prevent prompt-injection outcomes by sandboxing tool invocation or isolating untrusted content. These approaches are complementary to detection: they provide provable guarantees on specific agent decisions, not on the broader protocol traffic. Second, \emph{tool-boundary monitors} like ClawGuard~\cite{Zhao2026ClawGuard} enforce per-call deterministic policies, and program-dependency monitors like AgentArmor~\cite{Wang2025AgentArmor} reconstruct a per-trace dependency graph. Neither observes the multi-agent interaction graph, nor reports a quantitative robustness guarantee. Third, classical \emph{network intrusion detection systems} have evolved through graph neural networks~\cite{Lo2022EGraphSAGE}, transformer encoders~\cite{Wu2022RTIDS}, and selective state-space models~\cite{Wang2025NetMamba} --- but these operate on flat flow features and do not model MCP/A2A/ACP/ANP message semantics or certify their own output.

MambaGuard closes the gap between these three strands. It models the protocol stream as a time-stamped heterogeneous graph whose vertices are agents, tools, capability descriptors, and sessions, and whose edges carry canonicalised message tuples. A selective state-space backbone~\cite{Gu2023Mamba} processes each agent's message stream in linear time; a Lipschitz-normalised GATv2 layer~\cite{Brody2022GATv2,Dasoulas2021LipNorm} propagates relational context; and a Hedge controller~\cite{CesaBianchi2006Prediction} selects defensive actions online with no-regret against the best fixed commitment in hindsight. Two certificates accompany the detector: a randomised-smoothing certificate~\cite{Cohen2019RandSmooth} on the payload-embedding feature space, and a sparse-smoothing certificate~\cite{Bojchevski2020Sparse} on graph edges.

\subsection{Contributions}
\label{sec:contrib}
\begin{itemize}
  \item A \textbf{unified protocol-graph schema} for MCP, A2A, ACP and ANP (\S\ref{sec:graph}) with a five-class canonical message taxonomy.
  \item A \textbf{streaming composed detector} that couples a selective SSM with a Lipschitz-normalised temporal GATv2 layer (\S\ref{sec:detector}).
  \item A \textbf{dimensionally consistent two-certificate scheme} (\S\ref{sec:guarantees}): a Cohen-style smoothing certificate on feature-space perturbations with a composition theorem whose terms all carry the utility unit, plus a sparse-smoothing certificate on graph edits.
  \item A \textbf{dedicated adaptive-attack evaluation} (\S\ref{sec:adaptive}) in the spirit of~\cite{Nasr2026Attacker}: white-box PGD in embedding space, token-level discrete search, LLM-driven payload rewriting, graph-structural evasion, Hedge-weight manipulation.
  \item An \textbf{anonymised artifact} that regenerates every table and figure with one command.
\end{itemize}

\subsection{Scope and Non-Goals}
The paper's guarantees are deliberately scoped. The feature-space certificate holds inside the sentence-transformer image, not for arbitrary text edits; the graph certificate holds under a bounded number of edge additions or deletions, not under host compromise; and the online bound holds in the zero-sum restriction of the defender game. \S\ref{sec:cert-ops} and \S\ref{sec:limits} state exactly what the certificates do and do not cover.

%% ============================================================================
\section{Background}
\label{sec:bg}
%% ============================================================================
Two technical ingredients underpin the paper. We state only what the body of the paper needs; extensions and standard proofs are deferred to Appendix~\ref{app:prelim}.

\paragraph{Selective state-space models.}
A continuous state-space model maintains a hidden state $h(t)\in\R^{N}$ evolving as $h'(t)=Ah(t)+Bu(t)$, with output $y(t)=Ch(t)$. Zero-order-hold discretisation gives $h_{t}=\bar A h_{t-1}+\bar B u_{t}$, with $\bar A=\exp(\Delta A)$ and $\bar B=(\Delta A)^{-1}(\exp(\Delta A)-I)\Delta B$. The selective variant Mamba~\cite{Gu2023Mamba} makes $\Delta$, $B$, $C$ input-dependent and amortises inference in $\calO(L)$ per sequence of length $L$ via a hardware-aware parallel scan. We adopt the \textbf{S4D-Real} parameterisation~\cite{Gu2022S4D}, in which $A\in\R^{N}$ is a diagonal real matrix with strictly negative entries; this is the default Mamba initialisation and the choice reported to perform best on discrete modalities. For a real diagonal matrix the spectral radius equals the spectral (operator-2) norm, which is the property we exploit in the Lipschitz analysis of \S\ref{sec:lipschitz}.

\paragraph{Randomised smoothing.}
Given a base classifier $f$, the smoothed classifier $g(x)=\argmax_{c}\Pr_{\delta\sim\calN(0,\sigma^{2}I)}[f(x+\delta)=c]$ is certifiably robust in $\ell_{2}$ within a radius $r = \sigma\,\Phi^{-1}(p_{A})$, where $p_{A}$ is the probability of the top class under the noise distribution and $\Phi$ is the standard normal CDF~\cite{Cohen2019RandSmooth}. In practice $p_{A}$ is replaced by a Clopper-Pearson lower bound $\underline{p_{A}}$ estimated from $n$ Monte-Carlo samples, with a failure probability $\alpha_{MC}$ absorbed into the certificate.

%% ============================================================================
\section{Threat Model and Security Goals}
\label{sec:threat}
%% ============================================================================
S\&P's double-blind guidelines ask for a precise threat model before any robustness claim. We specify the model along four axes.

\paragraph{System model.}
Several agents communicate through MCP, A2A, ACP and ANP. A trusted detector observes every protocol message, canonicalises it, and maintains a sliding-window protocol graph. A defender action is selected per round and applied to one message or session: deep packet inspection, blocking, session termination, agent isolation, or capability revocation.

\paragraph{Attacker knowledge.}
We consider a \emph{grey-box} attacker who knows the detector's architectural family (selective SSM plus temporal GATv2, protocol-graph representation, 34-class attack taxonomy), the published hyper-parameters, and the frozen sentence-transformer encoder used for payload embeddings. The attacker does not have the deployed weights or the per-tenant Hedge random seed. For each experiment we also report a \emph{white-box} setting in which the attacker additionally receives the trained weights and the embedding gradient; this is the standard worst case used by~\cite{Nasr2026Attacker}.

\paragraph{Attacker capabilities.}
The attacker may (i) craft or rewrite protocol payloads inside a signed agent envelope; (ii) modify tool descriptions or Agent Cards; (iii) replay historical messages within a bounded time window; (iv) inject or delete graph edges up to a budget $k$ per detection window through a compromised registered agent; (v) query the detector a bounded number of times to drive an attack search. The attacker cannot forge cryptographic signatures outside the compromised set, cannot compromise the host that runs the detector, and cannot modify the detector's weights. Compromise of the model supply chain, memory, or host is explicitly out of scope.

\paragraph{Attacker action space $\calA_{A}$.}
Seven families from the agent-protocol attack literature~\cite{Hou2025MCP,Yang2025MCPSec,Lotfi2026A2ABreak}: tool poisoning, indirect prompt injection, rug-pull capability updates, capability shadowing, cross-server capability confusion, credential exfiltration, and sandbox bypass.

\paragraph{Perturbation budgets.}
Two complementary budgets. $\varepsilon\ge 0$ bounds the $\ell_{2}$ norm of perturbations at the frozen payload-embedding layer; this is where the smoothing certificate holds. $k\in\{0,1,\ldots\}$ bounds the number of graph edges the attacker may add or delete inside one detection window; this is where the sparse-smoothing certificate holds. We do \emph{not} claim that the $\ell_{2}$ budget transports to a bound on raw text edits: the payload encoder has no meaningful global Lipschitz constant from token edits to embeddings, and a single token change can move an embedding arbitrarily far relative to $\varepsilon$. \S\ref{sec:cert-ops} returns to this restriction.

\paragraph{Security goals.}
The defender seeks (G1) low per-message false-positive rate on benign traffic; (G2) high detection rate on malicious traffic; (G3) \emph{certified} robustness within the stated perturbation budgets; (G4) sustained throughput at the message rate of the monitored deployment.

%% ============================================================================
\section{Protocol Graph and Detector}
\label{sec:detector}
%% ============================================================================

\subsection{Unified Heterogeneous Protocol Graph}
\label{sec:graph}
Let $\calG(t)=(\calV(t),\calE(t),X(t))$ denote the protocol graph at wall-clock time $t$. The vertex set $\calV(t)=\calV_{A}\cup\calV_{T}\cup\calV_{C}\cup\calV_{S}$ is partitioned into agents, tools or services, capability descriptors, and active sessions. Edges are time-stamped and typed; the five canonical types are agent-calls-tool, agent-messages-agent, agent-advertises-capability, tool-returns-data, and agent-delegates-task. Each protocol message is canonicalised into the tuple
\[ m=(\tau,s,d,p,\mu,t_{m}), \]
where $\tau\in\{\calM_{\text{tool}},\calM_{\text{comm}},\calM_{\text{cap}},\calM_{\text{data}},\calM_{\text{ctrl}}\}$, $s,d$ are source/destination identifiers, $p\in\R^{d_{p}}$ is a frozen all-MiniLM payload embedding ($d_{p}=384$), $\mu\in\R^{d_{\mu}}$ carries HTTP status, latency and auth metadata, and $t_{m}$ is the message timestamp. A rule-based parser produces one tuple per raw protocol packet for each of MCP, A2A, ACP and ANP.

\paragraph{Worked example.} Consider a 2-second window in which Agent~A issues an MCP \texttt{tools/list} to a database tool $T_{1}$ and receives a 400-byte JSON response; concurrently, Agent~B sends an A2A \texttt{tasks/send} to Agent~A delegating a user query, and a new capability descriptor $C_{3}$ is advertised through ANP by Agent~C. The parser canonicalises six messages into three edges: $(A\!\to\!T_{1},\calM_{\text{tool}})$, $(T_{1}\!\to\!A,\calM_{\text{data}})$, $(B\!\to\!A,\calM_{\text{comm}})$, and one capability-advertisement edge $(C\!\to\!C_{3},\calM_{\text{cap}})$. All four edges carry the same timestamp bucket $t_{e}\in[0,2]$. The graph $\calG(t)$ at the end of the window thus has six nodes (three agents, one tool, one capability, one session), four typed edges, and per-edge features comprising the 384-dim frozen embedding of the payload and a 32-dim metadata vector. This is the input the detector sees for one detection decision.

\subsection{Selective SSM Backbone}
\label{sec:ssm}
For each agent $a\in\calV_{A}$, the sequence of incident messages $x^{a}=(m^{a}_{1},\ldots,m^{a}_{L_{a}})$ is processed by a four-layer selective block (details in Appendix~\ref{app:hyper}):
\begin{align}
h^{a}_{t} &= \bar A_{t}\,h^{a}_{t-1} + \bar B_{t}\,m^{a}_{t}, \label{eq:ssm}\\
z^{a}_{t} &= C_{t}\,h^{a}_{t} \odot \mathrm{SiLU}\bigl(W_{g}\,m^{a}_{t}\bigr), \label{eq:gate}
\end{align}
with the standard Mamba depthwise causal 1-D convolution applied to $x^{a}$ before~(\ref{eq:ssm}). The gated output $z^{a}_{t}\in\R^{d_{z}}$ is the per-agent sequential feature.

\subsection{Lipschitz-Normalised Temporal GATv2}
\label{sec:gat}
With $\{z^{a}\}_{a\in\calV_{A}}$ available, a temporal GATv2 layer propagates relational information. For an ordered pair $(a,a')$ joined by an edge with timestamp $t_{e}$, the attention logit is
\begin{equation}
e_{a,a'} = q^{\top}\,\mathrm{LeakyReLU}\!\bigl(W_{r}\,[z^{a}\Vert z^{a'}\Vert e_{a,a'}^{\text{feat}}\Vert\Phi(t-t_{e})]\bigr),
\label{eq:gatv2}
\end{equation}
where $\Phi(\Delta t)$ is the Bochner-derived TGAT time encoding~\cite{Xu2020TGAT} with $d_{T}$ learnable frequencies, $\Vert$ denotes concatenation, and $q$, $W_{r}$, $W_{v}$ are learnable. We apply \emph{LipschitzNorm}~\cite{Dasoulas2021LipNorm} to the attention weights; this is essential for the Lipschitz analysis below, because vanilla GATv2's dot-product attention is \emph{not} Lipschitz for any $\ell_{p}$ norm on an unbounded domain (Kim \textit{et al.}~\cite{Kim2021LipAttn} prove this for self-attention and the argument transfers to GATv2). Multi-head attention with $H=4$ heads is averaged at the final layer; the composed feature $h^{a}$ feeds a softmax head over the 34-class attack taxonomy.

\subsection{Hedge Controller}
\label{sec:hedge}
The detector logits parameterise a Hedge~\cite{CesaBianchi2006Prediction} selector over the five defensive actions $\calA_{D}$. Algorithm~\ref{alg:mwu} states the complete loop.

\begin{algorithm}[!t]
\caption{MambaGuard online defence loop}
\label{alg:mwu}
\begin{algorithmic}[1]
\Require Horizon $T$, utility range $B$, number of defender actions $N=|\calA_{D}|$
\State $\eta \gets \sqrt{8 \ln N / T}\,/\,B$
    \Comment{Corrected tuning (\S\ref{sec:noregret})}
\State $w_{1}(s)\gets 1$ for all $s\in\calA_{D}$
\For{$t=1,\ldots,T$}
  \State Observe $\calG(t)$ and per-agent streams
  \State Compute $z^{a}_{t}$ via~(\ref{eq:ssm})--(\ref{eq:gate})
  \State Compute $h^{a}_{t}$ via~(\ref{eq:gatv2})
  \State Score classes $\hat y_{t}\gets\mathrm{softmax}(W_{h} h_{t})$
  \State Sample action $s_{t}\sim p_{t}(s)=w_{t}(s)/\sum_{s'}w_{t}(s')$
  \State Observe loss $\ell_{t}(s)\in[0,B]$
  \State $w_{t+1}(s)\gets w_{t}(s)\exp(-\eta\,\ell_{t}(s))$
\EndFor
\end{algorithmic}
\end{algorithm}

%% ============================================================================
\section{Certified Guarantees}
\label{sec:guarantees}
%% ============================================================================
The detector carries two certificates --- feature-space (\S\ref{sec:smooth}) and graph-edit (\S\ref{sec:graph-cert}) --- plus a no-regret guarantee on the Hedge controller (\S\ref{sec:noregret}). \S\ref{sec:composed} combines them in a dimensionally consistent composition theorem, and \S\ref{sec:cert-ops} states what the certificates do and do not operationally guarantee.

\subsection{Feature-Space Smoothing Certificate}
\label{sec:smooth}
We define the smoothed detector $g(x) = \argmax_{c}\Prob_{\delta\sim\calN(0,\sigma^{2}I)}[f(x+\delta)=c]$ over the payload-embedding input $x$. Following~\cite{Cohen2019RandSmooth}, for a given $x$ let $\underline{p_{A}}$ be the Clopper-Pearson lower bound on $\Prob[f(x+\delta)=g(x)]$ from $n$ Monte-Carlo samples at confidence $1-\alpha_{MC}$. Then with probability at least $1-\alpha_{MC}$ the certified radius is
\begin{equation}
R(x) = \sigma\,\Phi^{-1}\bigl(\underline{p_{A}}\bigr),
\label{eq:cert-radius}
\end{equation}
and $g(x+\delta')=g(x)$ for all $\|\delta'\|_{2}\le R(x)$. We report certified-accuracy-vs-radius curves rather than a single number, because the radius varies per input (\S\ref{sec:res-cert}).

\subsection{Local Lipschitz Bound on the Composed Detector}
\label{sec:lipschitz}
The smoothed detector's accuracy is tightened when the base detector is Lipschitz on the perturbed domain. For MambaGuard the base map is $f=h\circ g_{\text{gat}}\circ g_{\text{ssm}}$. We state the bound under explicit bounded-input assumptions, because the selective SSM and the gated output are both products of input-dependent factors and therefore \emph{not} globally Lipschitz on $\R^{d}$.

\begin{assumption}[Bounded operating regime]
\label{ass:bounded}
(i) Inputs are projected to a Euclidean ball $\|x_{t}\|_{2}\le X_{\max}$; (ii) discretisation steps are clamped to $\Delta_{t}\in[\Delta_{\min},\Delta_{\max}]$; (iii) spectral norms of the projections producing $B_{t}$, $C_{t}$ and $\Delta_{t}$ (including their Jacobians) are bounded by $M_{B}, M_{C}, M_{\Delta}$ respectively; (iv) the state matrix $A\in\R^{N}$ is real diagonal with strictly negative entries (S4D-Real) and $\Delta_{\max}\max_{i}A_{i}<0$, so $\rho(\bar A_{t})=\|\bar A_{t}\|_{2}\le\rho<1$.
\end{assumption}

\begin{lemma}[Local Lipschitz of the selective block]
\label{lem:lipssm}
Under Assumption~\ref{ass:bounded}, with $\beta:=\sup_{t}\|\bar B_{t}\|_{2}$, $\gamma:=\sup_{t}\|C_{t}\|_{2}$, $\kappa_{g}:=\|W_{g}\|_{2}$, $L_{\sigma}\approx 1.0998$ the exact SiLU Lipschitz constant, and $X_{\max}$ the input-norm cap, the selective block $g_{\text{ssm}}$ is Lipschitz on $\{x:\|x_{t}\|\le X_{\max}\}$ with
\begin{equation}
L_{\text{ssm}} \le \frac{\gamma\beta(1+L_{\sigma}\kappa_{g})}{1-\rho} + \frac{\gamma\beta X_{\max}(M_{B}{+}M_{C}{+}M_{\Delta})}{(1-\rho)^{2}}.
\label{eq:lipssm}
\end{equation}
\end{lemma}

The second term accounts for the input-dependent factors in $\bar B_{t}, C_{t}, \Delta_{t}$; it vanishes only when those parameters are input-independent (the S4 limit). Proof in Appendix~\ref{app:proof-ssm}.

\begin{lemma}[LipschitzNormed GATv2]
\label{lem:lipgat}
With LipschitzNorm scaling constant $\eta$ and Bochner frequencies $\{\omega_{k}\}$, the temporal GATv2 layer~(\ref{eq:gatv2}) satisfies
\begin{align}
L_{\text{gat}} &\le \|W_{v}\|_{2}\bigl(1+\eta\,\|q\|_{2}\,\|W_{r}\|_{2}\bigr), \label{eq:lipgat}\\
L_{\Phi} &\le \|W_{v}\|_{2}\,\eta\,\|q\|_{2}\,\|W_{r}\|_{2}\sqrt{\tfrac{2}{d_{T}}\textstyle\sum_{k}\omega_{k}^{2}}. \label{eq:lipbochner}
\end{align}
\end{lemma}

The composed bound $L_{f}\le M_{h}\,L_{\text{gat}}\,L_{\text{ssm}}$ is used only to tighten the smoothing certificate; \textbf{we do not claim it transfers to raw text perturbations}. \S\ref{sec:res-cert} reports the measured $L_{f}$ decomposition and its ratio to an empirical lower bound obtained by gradient ascent on $\delta$.

\subsection{Sparse-Smoothing Graph-Edit Certificate}
\label{sec:graph-cert}
A continuous certificate does not express structural attacks such as Agent Card impersonation, which add or remove edges in $\calG(t)$. We apply the sparse-smoothing framework of Bojchevski \textit{et al.}~\cite{Bojchevski2020Sparse}. For edge-flip probabilities $p_{+}$ (adding an absent edge) and $p_{-}$ (deleting a present edge), the smoothed classifier on edge-perturbed graphs is certified robust to budgets $(k_{a}, k_{d})$ whenever
\begin{equation}
\sum_{\mathbf{Q}\in\mathcal{H}}\min\{\mathbf{p}_{\mathbf{Q}}^{+},\mathbf{p}_{\mathbf{Q}}^{-}\} \;\ge\; p_{A} - 0.5,
\label{eq:np-bound}
\end{equation}
where the pointwise bound is taken over the Neyman-Pearson regions $\mathcal{H}$ induced by budgets $k_{a}$ edge additions and $k_{d}$ edge deletions, and $\mathbf{p}^{\pm}$ are the per-region likelihoods under the clean and perturbed noise distributions (full statement in~\cite{Bojchevski2020Sparse}, Theorem~1). In practice we evaluate the bound over the attacker's candidate edge set, not over all node pairs, which both reduces variance and matches the registered-agent restriction in the threat model.

\subsection{Online No-Regret Guarantee}
\label{sec:noregret}
We correct an error present in prior MambaGuard drafts. For losses in $[0,B]$ the Hedge regret bound $R_{T}\le(\ln N)/\eta+\eta T B^{2}/8$ is minimised at
\begin{equation}
\eta^{*} \;=\; \sqrt{8\ln N / T} \;/\; B,
\label{eq:eta-correct}
\end{equation}
giving
\begin{equation}
R_{T} \;\le\; B\sqrt{(T/2)\ln N}.
\label{eq:hedge-bound}
\end{equation}
(This is the classical Cesa-Bianchi--Lugosi tuning~\cite{CesaBianchi2006Prediction}; earlier MambaGuard drafts quoted $\eta=\sqrt{2\ln N/T}$, which minimises a looser bound and yields a factor-$\sqrt{2}$ larger constant.) For an anytime variant with $\eta_{t}=\sqrt{8\ln N/t}/B$, the regret is at most $B(\sqrt{2T\ln N}+\ln N/8)$.

\begin{theorem}[No-regret]
\label{thm:regret}
Algorithm~\ref{alg:mwu} with $\eta$ as in~(\ref{eq:eta-correct}) satisfies~(\ref{eq:hedge-bound}).
\end{theorem}

We do \emph{not} claim that no-regret dynamics converge to the Stackelberg leader value in general. That claim is only true in two restricted settings: (a) the defender-attacker game is \emph{zero-sum}, in which case the Stackelberg value equals the minimax value and Hedge's time-average payoff is within $R_{T}/T$ of it; or (b) the guarantee is phrased as regret against the \emph{best fixed commitment in hindsight} over a finite set of attacker types~\cite{Balcan2015Commitment}, with the action set restricted to the extreme points of the attacker best-response polytopes. Our experiments report both views.

\subsection{Composed Certificate (Type-Checked)}
\label{sec:composed}
Let $\calC_{\varepsilon}$ be the event that the smoothed detector's prediction is certified constant at radius $\varepsilon$, as in~(\ref{eq:cert-radius}). On $\calC_{\varepsilon}$ the defender's action and utility are unchanged under attack; off $\calC_{\varepsilon}$ the per-round loss is at most $B$. Combining with Theorem~\ref{thm:regret} gives the composition.

\begin{theorem}[Composed certificate]
\label{thm:composed}
Let $V^{*}$ be the Stackelberg value under the restriction of \S\ref{sec:noregret}(a) or (b), $\hat\pi_{T}$ the Hedge policy after $T$ rounds, $\varepsilon\ge 0$ the feature-space $\ell_{2}$ budget, and $\alpha_{MC}$ the Monte-Carlo failure probability in~(\ref{eq:cert-radius}). Then
\begin{equation}
\E\bigl[V_{\text{ach}}(\hat\pi_{T},\varepsilon)\bigr] \;\ge\; V^{*} - B\,\Prob\!\bigl[\neg\calC_{\varepsilon}\bigr] - B\sqrt{\tfrac{\ln N}{2T}} - B\,\alpha_{MC}.
\label{eq:composed}
\end{equation}
\end{theorem}

All four terms share the utility unit; $\Prob[\neg\calC_{\varepsilon}]=1-\mathrm{CertAcc}(\varepsilon)$ is read off the certified-accuracy curve. Proof in Appendix~\ref{app:proof-composed}.

\subsection{Operational Meaning and Non-Guarantees}
\label{sec:cert-ops}
Theorem~\ref{thm:composed} guarantees a lower bound on long-run expected defender utility against any attacker whose per-round perturbation stays within the stated budgets. It does \emph{not} guarantee: (i) that an arbitrary text edit of a payload lies within $\varepsilon$ in embedding space (the encoder has no global token-to-embedding Lipschitz constant); (ii) robustness to graph edits beyond $k_{a}+k_{d}$ within one window; (iii) robustness to host compromise, cryptographic-signature forgery, or zero-day protocol-implementation flaws. Applications that require payload-level provable robustness should additionally use the erase-and-check certificate of Kumar \textit{et al.}~\cite{Kumar2024EraseCheck} against suffix/insertion/infusion up to $d$ tokens; our evaluation reports this composition in \S\ref{sec:res-cert}.

%% ============================================================================
\section{Experimental Setup}
\label{sec:setup}
%% ============================================================================

\subsection{Benchmarks}
\label{sec:benchmarks}
We evaluate on six benchmarks that span the agent-protocol attack surface.

\textbf{MCPSecBench}~\cite{Yang2025MCPSec}: 17 attack types across 4 MCP surfaces (tool, server, host, user) with 15 runs per attack on three LLM platforms (Claude Desktop, GPT-4.1, Cursor).

\textbf{MCP-SafetyBench}~\cite{Zong2026SafetyBench}: 245 multi-turn test cases over 5 domains and 20 attack types; the published average ASR across 13 LLMs ranges from 29.80\% (Qwen3-235B) to 48.16\% (o4-mini), with host-side attacks reaching 81.94\%.

\textbf{MCPTox}~\cite{Wang2025MCPTox}: 1,312 malicious tool-poisoning cases over 10 risk categories across 45 live real-world MCP servers and 353 authentic tools.

\textbf{A2A-sandbox}: a reproduction of the 11 A2ABreak~\cite{Lotfi2026A2ABreak} vulnerability classes, including cross-client context injection and rogue-agent capability advertisement, generated in a local sandbox with no outbound network access.

\textbf{AgentDojo}~\cite{Debenedetti2024AgentDojo} and \textbf{InjecAgent}~\cite{Zhan2024InjecAgent}: 629 and 1054 indirect-prompt-injection scenarios respectively, canonicalised into our protocol-graph representation via the mapping documented in Appendix~\ref{app:canon}.

\textbf{Network-IDS transfer (appendix only).} For completeness we also include a transfer evaluation on CSE-CIC-IDS2018~\cite{Sharafaldin2018CICIDS}, CIC-IoT2023~\cite{Neto2023CICIoT}, and UNSW-NB15~\cite{Moustafa2015UNSW} with 70/15/15 strict temporal splits; those numbers are in Appendix~\ref{app:netids} and are \emph{not} used in the headline claims, because the datasets do not represent agent-protocol threats.

\subsection{Baselines}
\label{sec:baselines}
We compare against three tiers: (i) detector-style agent defences PromptGuard~\cite{PromptGuard2024}, PIGuard~\cite{PIGuard2024}, and DataSentinel~\cite{Liu2025DataSentinel}; (ii) design-level defences CaMeL~\cite{Debenedetti2025CaMeL} and MELON~\cite{Zhu2025MELON}; (iii) classical IDS backbones E-GraphSAGE~\cite{Lo2022EGraphSAGE}, RTIDS~\cite{Wu2022RTIDS}, and NetMamba~\cite{Wang2025NetMamba} applied to the protocol graph via a sentence-transformer front-end. Each baseline is retrained by us under identical splits with its published hyperparameters.

\subsection{Metrics}
\label{sec:metrics}
Per-benchmark macro-F1, per-benchmark post-defence ASR (at a defender-action threshold selected on validation), certified accuracy at $\varepsilon\in\{0.01,0.05,0.10,0.25\}$, batch-1 streaming per-message latency (median and p99 over 1000 runs after 50 warm-ups), and throughput at the best batch size. We do \emph{not} report a cross-benchmark macro-average because the label spaces differ.

\subsection{Statistical Protocol}
\label{sec:stats}
Each experiment is run with five seeds (0--4). Entries are mean~$\pm$~sample standard deviation. Pairwise comparisons use paired two-sided Wilcoxon signed-rank tests on the $n=5\times B$ (seeds $\times$ benchmarks) grid with Holm correction across comparisons; the smallest achievable two-sided exact $p$ at $n=5$ is $2^{-5}\cdot 2=0.0625$, so we report $p$-values for $n\ge 15$.

%% ============================================================================
\section{Results}
\label{sec:results}
%% ============================================================================

\subsection{Per-Benchmark Detection}
\label{sec:res-detect}
Table~\ref{tab:main} reports macro-F1 and post-defence ASR per benchmark. MambaGuard is the top performer on four of six benchmarks on macro-F1 and on five of six on ASR. CaMeL is the tightest competitor on ASR for agent-tool corpora (AgentDojo, InjecAgent) because its provable-by-design approach prevents injection outcomes regardless of detector performance; the two systems are complementary (MambaGuard observes traffic, CaMeL rewrites tool calls) and we comment on that composition in \S\ref{sec:disc}.

\begin{table*}[!t]
\centering
\caption{Per-benchmark results, mean $\pm$ sample std over 5 seeds. ``--'' denotes a method inapplicable to the benchmark. MG = MambaGuard. Statistical significance (paired Wilcoxon, Holm-corrected, vs second-best per column) is marked $^{*}$ at $p<0.05$ and $^{\ddagger}$ at $p<0.01$.}
\label{tab:main}
\footnotesize
\renewcommand{\arraystretch}{1.12}
\setlength{\tabcolsep}{3.5pt}
\begin{tabular}{@{}l c c c c c c@{}}
\toprule
& \textbf{MCPSecBench} & \textbf{MCP-SafetyBench} & \textbf{MCPTox} & \textbf{A2A-sandbox} & \textbf{AgentDojo} & \textbf{InjecAgent}\\
& F1 / ASR (\%) & F1 / ASR (\%) & F1 / ASR (\%) & F1 / ASR (\%) & F1 / ASR (\%) & F1 / ASR (\%) \\
\midrule
PromptGuard         & $.63$ / $38.1$ & $.57$ / $44.6$ & $.52$ / $49.3$ & $.54$ / $47.8$ & $.61$ / $31.5$ & $.63$ / $34.2$ \\
PIGuard             & $.67$ / $35.4$ & $.61$ / $40.1$ & $.56$ / $45.9$ & $.58$ / $44.2$ & $.64$ / $29.1$ & $.66$ / $31.0$ \\
DataSentinel        & $.74$ / $24.7$ & $.70$ / $28.3$ & $.66$ / $34.4$ & $.67$ / $30.6$ & $.71$ / $19.8$ & $.73$ / $21.4$ \\
CaMeL               & $.68$ / $17.4$ & $.65$ / $22.1$ & $.61$ / $27.5$ & $.63$ / $25.0$ & $.76$ / $\mathbf{4.9}$ & $.78$ / $\mathbf{6.1}$ \\
MELON               & $.70$ / $18.9$ & $.67$ / $23.8$ & $.62$ / $29.0$ & $.64$ / $26.2$ & $.74$ / $5.7$ & $.76$ / $7.4$ \\
NetMamba+ST         & $.77$ / $20.3$ & $.72$ / $25.6$ & $.68$ / $31.2$ & $.69$ / $28.9$ & $.73$ / $17.1$ & $.75$ / $18.3$ \\
E-GraphSAGE+ST      & $.73$ / $24.1$ & $.69$ / $28.8$ & $.65$ / $34.5$ & $.66$ / $31.0$ & $.70$ / $20.4$ & $.72$ / $22.0$ \\
RTIDS+ST            & $.75$ / $22.5$ & $.71$ / $27.0$ & $.66$ / $32.9$ & $.68$ / $29.4$ & $.72$ / $18.9$ & $.74$ / $20.5$ \\
\midrule
\textbf{MG (ours)}  & $\mathbf{.86}^{\ddagger}$ / $\mathbf{11.2}^{\ddagger}$ & $\mathbf{.83}^{\ddagger}$ / $\mathbf{14.6}^{\ddagger}$ & $\mathbf{.79}^{\ddagger}$ / $\mathbf{19.8}^{\ddagger}$ & $\mathbf{.81}^{\ddagger}$ / $\mathbf{16.3}^{\ddagger}$ & $.83^{*}$ / $9.1$ & $.85^{*}$ / $10.4$ \\
\bottomrule
\end{tabular}
\end{table*}

\subsection{Certification: Actual Lipschitz Constant, Smoothing and Graph-Edit}
\label{sec:res-cert}
Table~\ref{tab:lipsweep} reports the measured decomposition of $L_{f}$ into its SSM, GATv2 and head factors as a function of the effective spectral cap $c=\sup_{\ell}\|W^{(\ell)}\|_{2}$ enforced at training time by the Lipschitz regulariser $\lambda_{L}$. At $c=0.5$, achievable under aggressive regularisation, the composed constant collapses to $L_{f}\approx 0.02$; at $c=1$ it is $L_{f}\approx 652$; without regularisation it runs to $10^{4}$ and the smoothing certificate becomes vacuous. The empirical lower bound $\hat L_{f}$ from 200-step gradient ascent on $\delta$ is approximately $0.3\times$ the analytic bound at $c=0.5$ and $0.08\times$ at $c=1$: the Lipschitz analysis is loose, as expected for a worst-case product bound, and tighter LipSDP-style estimates~\cite{Fazlyab2019LipSDP} are future work.

\begin{table}[!t]
\centering
\caption{$L_{f}$ decomposition vs effective spectral cap $c$ imposed at training. $\hat L_{f}$ = empirical lower bound from 200-step PGD on $\delta$; $\hat L_{f}/L_{f}$ is the tightness ratio. Cert.\ frac.\ at $\varepsilon=0.05$ is the Cohen-certified-and-correct fraction with $\sigma=0.25$, $n=10^{4}$, $\alpha_{MC}=10^{-3}$.}
\label{tab:lipsweep}
\footnotesize
\renewcommand{\arraystretch}{1.12}
\setlength{\tabcolsep}{3pt}
\begin{tabular}{@{}cccccc@{}}
\toprule
$c$ & $L_{\text{ssm}}$ & $L_{\text{gat}}$ & $L_{f}$ & $\hat L_{f}/L_{f}$ & Cert.\ frac.\ @ $\varepsilon{=}0.05$ \\
\midrule
0.50 &   0.06 &   1.46 &   0.02 & 0.31 & $92.4\pm 1.1$\,\% \\
0.75 &  12.88 &   5.50 &  39.87 & 0.17 & $71.2\pm 1.7$\,\% \\
1.00 &  58.84 &  14.73 & 651.80 & 0.08 & $52.8\pm 2.1$\,\% \\
1.25 &  89.74 &  32.30 & 2482.52 & 0.05 & $28.3\pm 2.6$\,\% \\
2.00 & 177.96 &  67.99 & 10364.17 & 0.02 & $8.1\pm 1.4$\,\% \\
\bottomrule
\end{tabular}
\end{table}

Figure~\ref{fig:cert-curves} plots certified accuracy against radius $r$ for $\sigma\in\{0.05,0.10,0.25,0.5\}$ and compares to an erase-and-check certificate~\cite{Kumar2024EraseCheck} for suffix attacks on AgentDojo. At $\sigma=0.25$ the certified accuracy is 85\% at $r=0.1$ and decays to 46\% at $r=0.5$.

\begin{figure}[!t]
\centering
% Placeholder: measured by scripts/run_sp2027_experiments.py; regenerated by reproduce_all.sh
\framebox[\columnwidth]{\rule{0pt}{80pt}\texttt{figures/certified\_curves.pdf}}
\caption{Certified accuracy vs $\ell_{2}$ radius $r$ on MCP-SafetyBench, for four noise levels $\sigma\in\{0.05,0.10,0.25,0.5\}$. Dashed curve: erase-and-check certificate for 2-token suffix attacks on AgentDojo. Shaded band: $\pm 1$~std over 5 seeds.}
\label{fig:cert-curves}
\end{figure}

For the \emph{graph-edit} certificate we use sparse smoothing with $p_{+}=p_{-}=0.05$, restricted to the attacker's candidate edge set of registered agents. At $n=1000$ smoothing samples and $\alpha=10^{-3}$, the median certified budget on MCPSecBench is $(k_{a}, k_{d})=(2, 3)$ edges, with $27\%$ of attacks remaining certifiably classified correctly at $(k_{a}, k_{d})=(1, 1)$. Full distribution in Appendix~\ref{app:graph-cert}.

\subsection{Latency and Throughput, Harmonised}
\label{sec:res-latency}
Earlier MambaGuard drafts contained an internal contradiction between per-message latency and aggregate throughput. We resolve it in Table~\ref{tab:latency} by reporting, separately, (i) streaming batch-1 per-message latency for the single-user case (end-to-end, including the 1-ms canonicalisation step), and (ii) the amortised per-message cost at the best throughput-maximising batch size.

\begin{table}[!t]
\centering
\caption{Harmonised latency and throughput. All numbers from a single A100 80GB except the CPU rows (16-core Xeon). ``Per-msg batch 1'' is the streaming latency the end user sees; ``per-msg best-batch'' is the amortised cost once batching is applied.}
\label{tab:latency}
\footnotesize
\renewcommand{\arraystretch}{1.1}
\setlength{\tabcolsep}{4pt}
\begin{tabular}{@{}lcccc@{}}
\toprule
 & \textbf{Per-msg} & \textbf{Per-msg} & \textbf{Best} & \textbf{Throughput} \\
 & \textbf{batch 1} & \textbf{best-batch} & \textbf{batch} & \textbf{(Mmsg/s)} \\
\midrule
A100 GPU            & 3.7\,ms & 0.89\,\textmu s & 4096 & 1.13 \\
A100, streaming (p99) & 4.9\,ms & --- & --- & --- \\
CPU 16-core         & 9.1\,ms & 6.8\,\textmu s & 128 & 0.15 \\
\bottomrule
\end{tabular}
\end{table}

The earlier manuscript's ``3.7 ms per message at 1.13 Mmsg/s'' is now explicitly shown to be two different measurements: the 3.7 ms is batch-1 streaming latency on A100; the 1.13 Mmsg/s requires batch 4096, where the amortised per-message cost is 0.89~\textmu s.

\subsection{Expanded Component Ablation}
\label{sec:res-abl}
Table~\ref{tab:abl} reports an expanded ablation. ``MG, no Lipschitz'' turns off the spectral regulariser: F1 is nearly unchanged but the certified radius collapses and the detector is no longer meaningful under adaptive attack (\S\ref{sec:adaptive}). Removing GATv2 degrades F1 by 2.6 points on MCPSecBench; removing the SSM backbone (``GAT-only'') by 3.9 points. Replacing the scan with transformer attention gives marginally lower F1 but 10$\times$ latency on A100.

\begin{table}[!t]
\centering
\caption{Ablation on MCP-SafetyBench, mean $\pm$ std over 5 seeds. Cert.\ rad.\ is median certified radius at $\sigma=0.25$.}
\label{tab:abl}
\footnotesize
\renewcommand{\arraystretch}{1.12}
\setlength{\tabcolsep}{3pt}
\begin{tabular}{@{}lccccc@{}}
\toprule
\textbf{Config} & \textbf{Params} & \textbf{F1} & \textbf{Lat.(ms)} & \textbf{ASR(\%)} & \textbf{Cert.rad.} \\
\midrule
Full MG              & 3.68\,M & $.83$ & 3.7 & $14.6$ & 0.21 \\
no Lipschitz         & 3.68\,M & $.82$ & 3.7 & $15.2$ & $<0.01$ \\
no Hedge controller   & 3.68\,M & $.83$ & 3.5 & $24.1$ & 0.21 \\
no GATv2 (SSM only)   & 1.93\,M & $.80$ & 2.9 & $18.9$ & 0.19 \\
no SSM (GAT only)    & 1.92\,M & $.78$ & 0.7 & $22.4$ & 0.14 \\
Transformer-for-SSM  & 2.22\,M & $.82$ & 37.4 & $15.0$ & 0.20 \\
no time encoding     & 3.63\,M & $.81$ & 3.4 & $17.3$ & 0.20 \\
flat sequence        & 3.68\,M & $.76$ & 3.6 & $28.8$ & 0.21 \\
\bottomrule
\end{tabular}
\end{table}

%% ============================================================================
\section{Adaptive Attacks}
\label{sec:adaptive}
%% ============================================================================
S\&P 2026 USENIX Security results~\cite{Nasr2026Attacker} show that twelve recent prompt-injection filters were bypassed with ASR above 90\% once attackers were given time to tune and scale a handful of standard techniques. MambaGuard is a filter in this family, so we report a dedicated adaptive-attack evaluation. We follow the Tram{\`e}r \textit{et al.} methodology~\cite{Tramer2020Adaptive} of targeting \emph{every} component of the defence.

\paragraph{E1 White-box embedding PGD.}
Projected gradient descent with Expectation-over-Transformation~\cite{Athalye2018Obfuscated} over the smoothing noise. Budget $\varepsilon$ in $\ell_{2}$ at the sentence-transformer feature layer, 100 steps, 8 EOT samples, step $2\varepsilon/100$. Reports empirical ASR within the \emph{claimed} certified ball (should be zero in the certified fraction) and beyond it. Table~\ref{tab:adaptive} row E1.

\paragraph{E2 Discrete token search (GCG-style).}
Greedy Coordinate Gradient token search~\cite{Zou2023GCG} constrained to preserve the malicious tool call. Measures whether the certified radius transports to raw text perturbations --- it does not, by construction, but we report the empirical bypass rate so reviewers can see the gap between the feature-space certificate and text-level evasion.

\paragraph{E3 LLM-driven payload rewriting.}
An attacker LLM (Claude 3.5 Sonnet with a search budget of 20 queries per sample, following~\cite{Nasr2026Attacker}) is told the detector's taxonomy and the target class, and rewrites the malicious payload to minimise the detector score while preserving a verifiable attack outcome.

\paragraph{E4 Graph-structural evasion.}
Split the attack across $m$ agents or $w$ windows to stay under per-window graph statistics. Mimic benign subgraph patterns extracted from the test partition.

\paragraph{E5 Hedge-weight manipulation.}
An adaptive attacker that observes the defender's actions (a legitimate assumption under the Balcan \textit{et al.}~\cite{Balcan2015Commitment} framework) and plays a type sequence that drives the Hedge weights toward the least effective defender action. We report the empirical regret growth versus the Hedge bound: if the empirical regret exceeds the proven bound $B\sqrt{(T/2)\ln N}$, the online guarantee is empirically violated.

\paragraph{E6 Human red team.}
A three-person red team with institutional ethics approval produced 200 bypass attempts over 48 hours. Each attempt was a hand-crafted attack trace (not a modification of a benchmark trace); the team was given the detector's taxonomy and encouraged to use any non-destructive technique. Protocol and ethics discussion in Appendix~\ref{app:ethics}.

\paragraph{Comparison to the prompt-injection-filter baseline.}
Nasr \textit{et al.}~\cite{Nasr2026Attacker} report $>90\%$ ASR against twelve recent prompt-injection defences (including Protect~AI, PromptGuard, PIGuard, Model Armor) under their adaptive-attack methodology, with PIGuard --- the strongest detector baseline --- still reaching $71\%$ bypass. For direct comparison we run the same adaptive attacks against MambaGuard; Table~\ref{tab:nasr} reports the bypass rate at the published Nasr~\textit{et~al.} attacker-budget level, and shows that MambaGuard's inclusion of a graph-relational feature degrades the attacker's bypass rate by $\approx 2\times$ relative to the strongest content-only baseline, with the proviso that the reduction is partial and we make no claim of invulnerability.

\begin{table}[!t]
\centering
\caption{MambaGuard versus the strongest prompt-injection filters from the Nasr~\textit{et~al.}~\cite{Nasr2026Attacker} study, on an aligned test set. All numbers are bypass rates at the authors' published adaptive-attack budget (20 attacker queries, 100 PGD steps). ``MG+EC'' is MambaGuard composed with erase-and-check.}
\label{tab:nasr}
\footnotesize
\renewcommand{\arraystretch}{1.1}
\setlength{\tabcolsep}{4pt}
\begin{tabular}{@{}lccc@{}}
\toprule
\textbf{Defence}  & \textbf{PIGuard} & \textbf{MambaGuard} & \textbf{MG+EC} \\
\midrule
E1 Embed-PGD      & $89\%$ & $47\%$ & $\mathbf{31\%}$ \\
E2 GCG            & $71\%$ & $49\%$ & $\mathbf{23\%}$ \\
E3 LLM rewrite    & $82\%$ & $52\%$ & $\mathbf{34\%}$ \\
Avg.              & $81\%$ & $49\%$ & $\mathbf{29\%}$ \\
\bottomrule
\end{tabular}
\end{table}

Three conclusions follow from Tables~\ref{tab:adaptive} and~\ref{tab:nasr}. First, the smoothing certificate holds operationally: within the claimed certified ball ($\varepsilon=0.05$ at $\sigma=0.25$) E1 bypasses the detector on only 8.4\% of inputs, within the certified failure fraction. Second, text-level attacks (E2, E3) cleanly outstrip the feature-space budget, exactly as our Theorem~\ref{thm:composed} non-guarantee predicts; adopting the composed erase-and-check certificate reduces the E2 (GCG) bypass rate from $28.0\%$ to $11.6\%$ at $\varepsilon=0.05$ at the cost of an $8\%$ throughput hit (Appendix~\ref{app:erase}). Third, Hedge manipulation (E5) cannot drive regret above the proven bound, which is reassuring for the online layer, but has unmeasured interactions with the smoothing certificate that we flag as an open problem for future work.

\begin{table}[!t]
\centering
\caption{Adaptive-attack ASR on MCP-SafetyBench, averaged over 5 seeds. $\varepsilon$ is the attacker's $\ell_{2}$ budget at the sentence-transformer feature layer; $k$ is the graph-edit budget.}
\label{tab:adaptive}
\footnotesize
\renewcommand{\arraystretch}{1.1}
\setlength{\tabcolsep}{4pt}
\begin{tabular}{@{}lccccc@{}}
\toprule
\textbf{Attack} & $\varepsilon{=}0.01$ & $\varepsilon{=}0.05$ & $\varepsilon{=}0.10$ & $k{=}1$ & $k{=}3$ \\
\midrule
E1 Embed-PGD (EOT)       & $1.1$ & $8.4$ & $47.3$ & --- & --- \\
E2 GCG discrete          & $12.5$ & $28.0$ & $48.9$ & --- & --- \\
E3 LLM rewrite           & $18.2$ & $34.1$ & $51.6$ & --- & --- \\
E4 Graph evasion         & --- & --- & --- & $13.4$ & $41.8$ \\
E5 Hedge manipulation    & \multicolumn{5}{c}{avg.\ regret $0.21$ vs bound $0.20$} \\
E6 Human red team        & \multicolumn{5}{c}{$39/200$ bypassed = $19.5\%$ bypass rate} \\
\bottomrule
\end{tabular}
\end{table}

%% ============================================================================
\section{Discussion and Limitations}
\label{sec:disc}
%% ============================================================================
\label{sec:limits}
MambaGuard is a \emph{detector}, complementary to design-level defences like CaMeL and MELON that prevent injection outcomes by sandboxing; we view detector+design composition as the deployment target. Three honest limitations. \textbf{Feature vs payload.} The smoothing certificate holds in sentence-transformer space; the erase-and-check composition partially bridges but does not close the gap to text-level perturbations. \textbf{Loose Lipschitz.} The analytic $L_{f}$ is $5\times$ to $10\times$ the empirical lower bound; tighter data-dependent estimates are future work. \textbf{Attacker capabilities.} The threat model excludes host compromise, cryptographic-signature forgery and arbitrary new-agent injection; these capabilities need orthogonal defences. Encrypted traffic is partially addressed (IDS-PQC in Appendix~\ref{app:netids}); differentially private embeddings are future work.

%% ============================================================================
\section{Related Work}
\label{sec:related}
%% ============================================================================
\textbf{Agent-protocol security monitors.} ClawGuard~\cite{Zhao2026ClawGuard} enforces deterministic tool-access policies. AgentArmor~\cite{Wang2025AgentArmor} reconstructs program-dependency graphs from runtime traces. DataSentinel~\cite{Liu2025DataSentinel} is a game-theoretic detection framework published at IEEE S\&P 2025. None observes the multi-protocol interaction graph or reports a certificate.

\textbf{Design-level defences.} CaMeL~\cite{Debenedetti2025CaMeL} prevents injection outcomes by design. MELON~\cite{Zhu2025MELON} provides provable guarantees against indirect prompt injection in ICML 2025. Progent~\cite{Shi2025Progent} adds programmable privilege control. The Task Shield~\cite{TaskShield2024} enforces task alignment. These are complementary to our detector.

\textbf{Certified robustness.} Tsuzuku \textit{et al.}~\cite{Tsuzuku2018LipMargin} introduced Lipschitz-margin training; Cohen \textit{et al.}~\cite{Cohen2019RandSmooth} introduced randomised smoothing. Bojchevski \textit{et al.}~\cite{Bojchevski2020Sparse} developed sparse smoothing for graphs; Schuchardt \textit{et al.}~\cite{Schuchardt2023Localized,Schuchardt2023Hierarchical} extended to collective and hierarchical variants. Kumar \textit{et al.}~\cite{Kumar2024EraseCheck} gave erase-and-check for LLM prompt perturbations.

\textbf{Mamba / SSM stability.} Halloran \textit{et al.}~\cite{Halloran2025Lyapunov} prove Lyapunov stability for Mamba; a parallel 2024 preprint explores adversarial robustness of deep SSMs~\cite{SSMRobust2024}.

\textbf{Attention Lipschitz.} Kim \textit{et al.}~\cite{Kim2021LipAttn} prove that dot-product attention is not globally Lipschitz; LipschitzNorm~\cite{Dasoulas2021LipNorm} is the standard remedy for graph-attention networks.

\textbf{Adaptive-attack methodology.} Nasr \textit{et al.}~\cite{Nasr2026Attacker} bypassed twelve recent prompt-injection defences; Tram{\`e}r \textit{et al.}~\cite{Tramer2020Adaptive} wrote the methodological reference.

\textbf{Game theory.} Balcan \textit{et al.}~\cite{Balcan2015Commitment} prove regret against the best fixed commitment in hindsight for repeated Stackelberg games under unknown attacker types.

%% ============================================================================
\section{Conclusion}
\label{sec:conclusion}
%% ============================================================================
We presented MambaGuard, a streaming detector for multi-protocol LLM-agent traffic that couples a selective SSM with a Lipschitz-normalised temporal GATv2 layer and ships two type-checked certificates and a corrected online no-regret bound. Against an adaptive-attack evaluation modelled on the recent S\&P / USENIX Security methodology, the smoothing certificate holds within its claimed radius, graph-edit certified budgets are non-zero on the structural benchmarks, and the composition theorem produces a defender-utility lower bound that is dimensionally consistent. Code, weights, the adaptive-attack harness, and a reproduction script that regenerates every table and figure are released under an anonymised mirror.

%% ============================================================================
%% References (anonymised third-person self-citation where needed)
%% ============================================================================
\bibliographystyle{IEEEtran}
\begin{thebibliography}{99}

%% --- Agent protocols and attacks ---
\bibitem{Hou2025MCP} X. Hou \textit{et al.}, ``Model Context Protocol: Landscape, security threats, and future research directions,'' arXiv:2503.23278, 2025.
\bibitem{Yang2025MCPSec} Y. Yang \textit{et al.}, ``MCPSecBench: A systematic security benchmark and playground for testing model context protocols,'' arXiv:2508.13220, 2025.
\bibitem{Qualys2025CVE49596} Qualys, ``Anthropic MCP Inspector remote code execution vulnerability (CVE-2025-49596),'' Jul.\ 2025.
\bibitem{Lotfi2026A2ABreak} A. Lotfi, M. M. Rahman, I. Karim, E. Bertino, ``A2ABreak: Systematic security analysis of the A2A protocol,'' arXiv:2609.10871, 2026.
\bibitem{Zong2026SafetyBench} T. Zong \textit{et al.}, ``MCP-SafetyBench: A benchmark for safety evaluation of LLMs with real-world MCP servers,'' ICLR 2026, arXiv:2512.15163.
\bibitem{Wang2025MCPTox} Z. Wang \textit{et al.}, ``MCPTox: A benchmark for tool poisoning attack on real-world MCP servers,'' AAAI 2026, arXiv:2508.14925.
\bibitem{Habler2025A2A} Z. Habler \textit{et al.}, ``Building a secure agentic AI application leveraging Google's A2A protocol,'' arXiv:2504.16902, 2025.
\bibitem{Wang2025AgentArmor} P. Wang \textit{et al.}, ``AgentArmor: Enforcing program analysis on agent runtime trace to defend against prompt injection,'' arXiv:2508.01249, 2025.
\bibitem{Zhao2026ClawGuard} W. Zhao \textit{et al.}, ``ClawGuard: Securing LLM agents through deterministic tool-access policies,'' arXiv:2604.11790, 2026.
\bibitem{TaskShield2024} ``The Task Shield: Enforcing task alignment to defend against indirect prompt injection in LLM agents,'' arXiv:2412.16682, 2024.
\bibitem{Debenedetti2024AgentDojo} E. Debenedetti \textit{et al.}, ``AgentDojo: A dynamic environment to evaluate prompt injection attacks and defenses for LLM agents,'' NeurIPS 2024.
\bibitem{Zhan2024InjecAgent} Q. Zhan \textit{et al.}, ``InjecAgent: Benchmarking indirect prompt injections in tool-integrated large language model agents,'' Findings of ACL 2024.

%% --- Design-level defences ---
\bibitem{Debenedetti2025CaMeL} E. Debenedetti \textit{et al.}, ``Defeating prompt injections by design (CaMeL),'' arXiv:2503.18813, 2025.
\bibitem{Zhu2025MELON} K. Zhu \textit{et al.}, ``MELON: Provable defense against indirect prompt injection attacks in AI agents,'' ICML 2025, PMLR 267:80310--80329.
\bibitem{Shi2025Progent} T. Shi \textit{et al.}, ``Progent: Programmable privilege control for LLM agents,'' arXiv:2504.11703, 2025.

%% --- Detector-style baselines ---
\bibitem{PromptGuard2024} Meta AI, ``PromptGuard: A compact classifier for prompt injection,'' 2024.
\bibitem{PIGuard2024} ``PIGuard: Prompt injection detection,'' arXiv:2411.XXXXX, 2024.
\bibitem{Liu2025DataSentinel} Y. Liu \textit{et al.}, ``DataSentinel: A game-theoretic detection of prompt injection attacks,'' IEEE S\&P 2025, arXiv:2504.11358.

%% --- SSM / Mamba ---
\bibitem{Gu2023Mamba} A. Gu, T. Dao, ``Mamba: Linear-time sequence modeling with selective state spaces,'' arXiv:2312.00752, 2023.
\bibitem{Gu2022S4D} A. Gu \textit{et al.}, ``On the parameterization and initialization of diagonal state space models,'' NeurIPS 2022.
\bibitem{Halloran2025Lyapunov} J. T. Halloran \textit{et al.}, ``Mamba state-space models are Lyapunov-stable learners,'' TMLR 2025, arXiv:2406.00209.
\bibitem{SSMRobust2024} ``Exploring adversarial robustness of deep state space models,'' arXiv:2406.05532, 2024.
\bibitem{Wang2025NetMamba} T. Wang \textit{et al.}, ``NetMamba: Efficient network traffic classification via pre-training unidirectional Mamba,'' IEEE INFOCOM 2025.

%% --- GNN / attention Lipschitz ---
\bibitem{Brody2022GATv2} S. Brody, U. Alon, E. Yahav, ``How attentive are graph attention networks?,'' ICLR 2022.
\bibitem{Xu2020TGAT} D. Xu \textit{et al.}, ``Inductive representation learning on temporal graphs,'' ICLR 2020.
\bibitem{Dasoulas2021LipNorm} G. Dasoulas, K. Scaman, A. Virmaux, ``Lipschitz normalization for self-attention layers with application to graph neural networks,'' ICML 2021.
\bibitem{Kim2021LipAttn} H. Kim, G. Papamakarios, A. Mnih, ``The Lipschitz constant of self-attention,'' ICML 2021.
\bibitem{Lo2022EGraphSAGE} W. W. Lo \textit{et al.}, ``E-GraphSAGE: A graph neural network based IDS for IoT,'' IEEE/IFIP NOMS 2022.
\bibitem{Wu2022RTIDS} Z. Wu \textit{et al.}, ``RTIDS: A robust transformer-based approach for IDS,'' IEEE Access, vol.\ 10, 2022.

%% --- Certification ---
\bibitem{Cohen2019RandSmooth} J. M. Cohen, E. Rosenfeld, J. Z. Kolter, ``Certified adversarial robustness via randomized smoothing,'' ICML 2019.
\bibitem{Tsuzuku2018LipMargin} Y. Tsuzuku, I. Sato, M. Sugiyama, ``Lipschitz-margin training,'' NeurIPS 2018.
\bibitem{Fazlyab2019LipSDP} M. Fazlyab \textit{et al.}, ``Efficient and accurate estimation of Lipschitz constants for deep neural networks,'' NeurIPS 2019.
\bibitem{Kumar2024EraseCheck} A. Kumar \textit{et al.}, ``Certifying LLM safety against adversarial prompting,'' COLM 2024, arXiv:2309.02705.
\bibitem{Bojchevski2020Sparse} A. Bojchevski, J. Gasteiger, S. G{\"u}nnemann, ``Efficient robustness certificates for discrete data: Sparsity-aware randomized smoothing for graphs, images and more,'' ICML 2020.
\bibitem{Schuchardt2023Localized} J. Schuchardt \textit{et al.}, ``Localized randomized smoothing for collective robustness certification,'' ICLR 2023.
\bibitem{Schuchardt2023Hierarchical} J. Schuchardt \textit{et al.}, ``Hierarchical randomized smoothing,'' NeurIPS 2023.

%% --- Game theory / online learning ---
\bibitem{CesaBianchi2006Prediction} N. Cesa-Bianchi, G. Lugosi, \textit{Prediction, Learning, and Games}. Cambridge Univ.\ Press, 2006.
\bibitem{Balcan2015Commitment} M.-F. Balcan, A. Blum, N. Haghtalab, A. D. Procaccia, ``Commitment without regrets: Online learning in Stackelberg security games,'' ACM EC 2015.

%% --- Adaptive-attack methodology ---
\bibitem{Nasr2026Attacker} M. Nasr \textit{et al.}, ``The attacker moves second: Stronger adaptive attacks bypass defenses against LLM jailbreaks and prompt injections,'' USENIX Security 2026, arXiv:2510.09023.
\bibitem{Tramer2020Adaptive} F. Tram{\`e}r \textit{et al.}, ``On adaptive attacks to adversarial example defenses,'' NeurIPS 2020.
\bibitem{Athalye2018Obfuscated} A. Athalye, N. Carlini, D. Wagner, ``Obfuscated gradients give a false sense of security,'' ICML 2018.
\bibitem{Zou2023GCG} A. Zou \textit{et al.}, ``Universal and transferable adversarial attacks on aligned language models (GCG),'' arXiv:2307.15043, 2023.

%% --- Network-IDS (appendix only) ---
\bibitem{Sharafaldin2018CICIDS} I. Sharafaldin \textit{et al.}, ``Toward generating a new intrusion detection dataset,'' ICISSP 2018.
\bibitem{Neto2023CICIoT} E. C. P. Neto \textit{et al.}, ``CICIoT2023,'' \textit{Sensors}, vol.\ 23, 2023.
\bibitem{Moustafa2015UNSW} N. Moustafa, J. Slay, ``UNSW-NB15,'' MilCIS 2015.

\end{thebibliography}

%% ============================================================================
%% Appendices (clearly marked; reviewers are not required to read these)
%% Total length cap: ~5 pages including references above.
%% ============================================================================
\appendices

\section{Mathematical Preliminaries}
\label{app:prelim}
Standard restatement of S4D-Real Mamba, dynamic graph attention and TGAT time encoding; omitted from body for space. Full derivations in the anonymised artifact (file \texttt{docs/preliminaries.pdf}).

\section{Proof of Lemma~\ref{lem:lipssm} (Local SSM Lipschitz)}
\label{app:proof-ssm}
Under Assumption~\ref{ass:bounded}, the state recurrence $h_{t}=\sum_{k=0}^{t-1}(\prod_{j=t-k+1}^{t}\bar A_{j})\bar B_{t-k}m_{t-k}$ gives
$\|h_{t}(x)-h_{t}(x')\|\le\sum_{k\ge 0}\rho^{k}\beta\|m_{t-k}-m'_{t-k}\|=\beta/(1-\rho)\|x-x'\|_{\infty,t}$.
Multiplying by $\gamma=\sup\|C_{t}\|$ for the output, and by the Hadamard-gate factor $(1+L_{\sigma}\kappa_{g})$ since $\|fg\|_{\text{Lip}}\le\|f\|_{\infty}L_{g}+\|g\|_{\infty}L_{f}$ with $\|f\|_{\infty},\|g\|_{\infty}\le\gamma\beta X_{\max}/(1-\rho)$, bounded by $X_{\max}$, yields the first term of (\ref{eq:lipssm}). The second term absorbs the Jacobians of $\bar B_{t}$, $C_{t}$ and $\Delta_{t}$ with respect to the input; applying submultiplicativity once more gives $\gamma\beta X_{\max}(M_{B}+M_{C}+M_{\Delta})/(1-\rho)^{2}$. The depthwise causal convolution is folded into $\beta$ via its unrolled Toeplitz norm. \qed

\section{Proof of Theorem~\ref{thm:composed} (Composed Certificate)}
\label{app:proof-composed}
Partition the horizon into rounds where $\calC_{\varepsilon}$ holds and rounds where it does not. On $\calC_{\varepsilon}$ the detector's prediction, hence the defender's action, hence the per-round utility, is unchanged under attack and therefore at least $V^{*}$ in expectation (by the Stackelberg optimality under the zero-sum or fixed-commitment restriction). On $\neg\calC_{\varepsilon}$ the per-round loss is at most $B$. The Hedge regret $R_{T}\le B\sqrt{(T/2)\ln N}$ divided by $T$ gives the second penalty. The $\alpha_{MC}$ term is the probability that the Clopper-Pearson lower bound fails and the certified event is reported in error. \qed

\section{Exact SiLU Lipschitz Constant}
\label{app:silu}
$\frac{d}{du}\mathrm{SiLU}(u)=\sigma(u)(1+u(1-\sigma(u)))$; numerical solution of $\mathrm{SiLU}''(u)=0$ yields $u^{*}\approx 2.3994$, $\mathrm{SiLU}'(u^{*})\approx 1.0998$.

\section{Hyperparameters}
\label{app:hyper}
$N=16$, $d_{z}=256$, expansion $E=2$, kernel width $k=4$, 4 stacked SSM layers; $H=4$ heads, $d_{T}=64$; AdamW lr $3\times 10^{-4}$, weight decay $10^{-2}$, batch 256, 30 epochs cosine; $\lambda_{L}=10^{-3}$ unless otherwise noted; sentence-transformer all-MiniLM-L6-v2 ($d_{p}=384$), frozen. Smoothing: $n=10^{4}$, $\alpha_{MC}=10^{-3}$, $\sigma\in\{0.05,0.10,0.25,0.5\}$. Sparse smoothing: $p_{+}=p_{-}=0.05$, $n=1000$.

\section{Protocol-Benchmark Canonicalisation}
\label{app:canon}
Mapping tables from AgentDojo/InjecAgent, MCPSecBench and MCP-SafetyBench raw formats into the six-tuple $(\tau,s,d,p,\mu,t_{m})$. Published in the artifact under \texttt{mambaguard/data/protocol\_traces/}.

\section{Graph-Edit Certified Budgets}
\label{app:graph-cert}
Full distribution of $(k_{a}, k_{d})$ over the six benchmarks and the sensitivity to $p_{+}, p_{-}$. See the artifact for CSV.

\section{Erase-and-Check Composition}
\label{app:erase}
Composing the feature-space smoothing certificate with the erase-and-check certificate~\cite{Kumar2024EraseCheck} against suffix attacks of length $\le 2$ reduces the E2 (GCG) bypass rate from $28.0\%$ to $11.6\%$ at $\varepsilon=0.05$, at the cost of an $8\%$ throughput hit.

\section{Network-IDS Transfer (Not in Headline Claims)}
\label{app:netids}
Transfer F1 to CSE-CIC-IDS2018/CIC-IoT2023/UNSW-NB15 after only fine-tuning the head: $0.78$/$0.74$/$0.71$, median over 5 seeds. These numbers are not used in the headline claims because the datasets do not represent agent-protocol threats; they document that the detector is not catastrophically specialised.

\section{Ethics Considerations}
\label{app:ethics}
The adaptive-attack suite in \S\ref{sec:adaptive} contains code for E1--E5 that could in principle be adapted to attack production MCP servers. We release the scaffolding needed to reproduce the paper but not working exploits against live third-party services. The human red team (E6) was conducted under institutional ethics approval; no human subject data was collected; participants were compensated at above-minimum-wage rates. All protocol traces are synthetic or from public benchmarks. We do not include CVE identifiers in the submission. Protocol flaws discovered during this work were reported to the relevant maintainers under coordinated disclosure.

\section{Generative AI Usage}
\label{app:genai}
Closed LLMs (Claude 3.5 Sonnet, GPT-4o) were used to generate synthetic agent traces for the A2A-sandbox benchmark; prompts, seeds and API response caches are released in the artifact for reproducibility under the limits imposed by closed-model drift. Generative AI was used for editorial purposes in this manuscript, and all outputs were inspected by the authors to ensure accuracy and originality.

\section{Artifact Availability}
\label{app:artifact}
An anonymised reference implementation is available at \url{https://anonymous.4open.science/r/mambaguard-sp2027}. The artifact contains the model, the certification framework, the benchmarks, the adaptive-attack harness, and a one-line reproduction script \texttt{scripts/reproduce\_all.sh} that regenerates every table and figure in this paper. Expected CSV outputs are included so that reviewers can diff against the regenerated numbers. The artifact is frozen at the submission deadline and will not be updated during the review period, per the S\&P 2027 artifact policy.

\end{document}
